%!TEX root = ../thesis.tex
% ************************** Thesis Abstract *****************************
% 4eme de couverture
\ifthispageodd{\newpage\thispagestyle{empty}\null\newpage}{}
\thispagestyle{empty}
\setlength{\columnseprule}{0pt}
\setlength\columnsep{10pt}
\newgeometry{
left=16mm,
top=18mm,
right=16mm,
bottom=20mm
}
%\fontfamily{rm}\selectfont

\lhead{}
\rhead{}
\rfoot{}
\cfoot{}
\lfoot{}

\noindent 
%*****************************************************
%***** LOGO DE L'ED À CHANGER ÉVENTUELLEMENT *********
%*****************************************************
\begin{textblock*}{61mm}(12mm,0mm)
	\noindent\includegraphics[height=24mm]{Figs/EOBE.png} %%ed%%
\end{textblock*}
%*****************************************************

%\newgeometry{margin=0.5in}
\begin{mdframed}[linecolor=Prune,linewidth=1]
\fontsize{7}{9}\selectfont
\vspace{-.25cm}
\paragraph*{Titre:} Analyse d’une nouvelle topologie fiable de convertisseur analogique-numérique pour l'environnement automobile

\begin{small}
\vspace{-.5cm}
\paragraph*{Mot clés:} 
\fontsize{7}{9}\selectfont
CAN, Low-OSR, électronique, automobile, température
\vspace{-.5cm}
\begin{multicols}{2}
\fontsize{7}{9}\selectfont
\paragraph*{Résumé:} 
La tendance du secteur automobile à développer des capteurs et actionneurs intelligents, faire cohabiter l'électronique analogique et l'électronique numérique devient un art. Placé au sein des actionneurs, pour la sécurité et le confort des passagers, les convertisseurs analogique-numérique (CAN) sont les composants clés de ces systèmes intelligents. Un CAN rapide, précis, et peu cher se révèle être un précieux allié pour les équipementiers automobiles. Pour diminuer les coûts, et faciliter l’utilisation de ce bloc, la surface de silicium occupée doit être considérablement réduite à moins de 0.5$\rm mm^2$. Quant à la précision du convertisseur, 12-bits tous les 5 coups d'une horloge de 100 MHz sont nécessaires pour une température de -40\(\degree \)C à 175\(\degree \)C.

Ce travail de recherche se focalise sur l'amélioration de l'efficacité énergétique sous les contraintes que l'environnement automobile représente. Notre principale contribution réside dans le développement par une approche top-down d'une nouvelle architecture à 3 étages de topologies différentes. Le premier étage est un $\Sigma\Delta$-Incrémental intrinsèquement linéaire. Le second étage est un algorithmique pour augmenter rapidement la résolution. Enfin, un SAR accroît la résolution avec faible consommation de puissance et surface de silicium.

Suite à l'analyse de 40 années d'état de l'art, la nouvelle architecture proposée fut validée par vérification des non-linéarités statiques (DNL, INL) à différents niveaux de modélisation. Commençant par un modèle MATLAB sans les limitations analogiques, le niveau de modélisation se raffine petit à petit jusqu'au niveau transistor du convertisseur. Un modèle Verilog-A permit la détermination des spécifications minimales des briques de base analogiques: les comparateurs et les amplificateurs à transconductance. La sensibilité de ces derniers à la température fut analysée pour limiter les erreurs commises sur les tensions analogiques. Une fois dessinés et les parasites extraits, les modèles variant avec la température remplacent leur modèle Verilog-A respectif afin d'obtenir les performances finales. Parallèlement, deux architectures de comparateurs ont été évaluées en température au sein d'une première puce de test. Deux méthodes ont été utilisées pour estimer l'offset des comparateurs, et un nouveau circuit asynchrone estime le délai. Une seconde puce de test permet de vérifier la sensibilité du SAR à la température malgré un fonctionnement pseudo-asynchrone.

Pour les comparateurs, le nouveau circuit de mesure différentielle du retard montre une précision de 60 ps dans le pire des cas, pour la plus petite surface sur puce connue en considérant la technologie utilisée. Comme la variation du retard est dépendante de la température, le choix d'un Strong-ARM (SA) ou d’un Double-Tail (DT) dépendra du bruit, de la puissance, de la tension d'alimentation, et de la spécification de kickback. Pour une tension d'alimentation standard, les SA comparateurs ciblent les systèmes à faible consommation avec une tolérance élevée pour le kickback différentiel. Au contraire, les DT comparateurs acceptent une plage de tension d'alimentation plus faible, et présentent un faible kickback différentiel, mais un bruit plus important. Testé de -40\(\degree \)C à 200\(\degree \)C, le dernier étage du CAN proposé, n'a pas besoin d'être calibré jusqu'à 180\(\degree \)C. Les résultats encourageants sur cet étage permettent la réutilisation de celui-ci pour calibrer les étages précédents. Et pour le CAN, nous estimons une résolution possible de 11,2 bits en 5 cycles d'horloge par échantillon avec une extension à 13,3 bits en 6 cycles d'horloge. La surface estimée est de 0,12$\rm mm^2$.

La puce de test pour le CAN est en cours de finalisation, une première étape sera sa caractérisation. Les résultats de cette session de mesure détermineront s'il est possible de pousser l'architecture à des fréquences plus élevées pour ensuite tirer parti du traitement numérique pour conserver les performances.
\end{multicols}
\end{small}
\end{mdframed}
\vspace*{-5mm}
\begin{mdframed}[linecolor=Prune,linewidth=1]
\fontsize{7}{9}\selectfont
\vspace{-.25cm}
\paragraph*{Title:} A New ADC topology for reliable conversion in the automotive environment

\begin{small}
\vspace{-.5cm}
\paragraph*{Keywords:} 
\fontsize{7}{9}\selectfont
ADC, Low-OSR, electronic, automotive, temperature
\vspace{-.5cm}
\begin{multicols}{2}
\fontsize{7}{9}\selectfont
\paragraph*{Abstract:} 
%-- paragraph de presentation
In the automotive industry, the trend being to develop smart sensors and actuators, the on-board electronic has been ever more an artful work to combine analog electronics and the digital one. While many monitoring and control systems play a crucial role as well for the safety as for the comfort of passengers, small components, like ADCs, are mandatory as a building block or as an essential functionality integrated into smart actuators. To that extent, a low-cost, fast and accurate analog to digital converter operating in those harsh conditions is a good ally for equipment manufacturers. To decrease the cost, the area is of primary concern. Considering re-use of the ADC as an IP-bloc, the area has been limited to less than half a square millimeter for an low-oversampling ratio of 5 to output a 12-bit code at a sample rate of 20 MSamples/s, over a wide temperature range from -40\(\degree \)C to 175\(\degree \)C.

%-- paragraph sur le but
This work focuses on the design of high-precision, high-speed and energy efficient ADC under the harsh environment the automotive one represents. Our main contribution relies on the development of an new hybrid topology proposal using 3 stages to cope with such constraints based on a top-down approach: A first counting stage inherently linear, an algorithmic stage allowing to increase rapidly the precision, and a SAR stage, ideal in terms of area and consumption, for a low number of bits.

%-- paragraph sur la method
Based on a 40 years literature review, a new topology proposal has been validated by checking its static metric of non-linearity (DNL, INL) at different level of modelisation. Starting by a MATLAB implementation without analog limitations, we refined step by step the model till we reach a transistor level of the ADC\@. Thence, Verilog-A model allows us to fix the minimum requirements of the key analog building blocks, to wit comparators and OTA\@. The latter has been analysed in order to limit the settling error sensitivity to the temperature. Laid-out, parasitic extracted simulation results of these considering PVT variations, they replace then previous high-level model to give final performances. Meanwhile, two well-known comparator architectures have been assessed as IP blocs inside a first test chip. To perform the offset extraction, both a conventional and a feedback loop have been inspected. To assess, the delay a new asynchronous circuit has been proposed. A second chip tests the sensitivity of the SAR to validate both the pseudo-asynchronous digital scheme, and a Double-Tail comparator in real operating conditions.

%-- paragraph sur les resultats
For comparators, the new differential measurement circuit of the delay demonstrate an accuracy of 60 ps in the worst case, over a large temperature range for the smallest chip area known with respect to the technology node size. The temperature variation of the delay being temperature dependent, the choice of a Strong-ARM or a Double-Tail hinge on the noise, power, supply voltage, and kickback specification. For standard power supply voltage, the Strong-ARM latch targets low-power systems application with a high tolerance for differential kickback. To the contrary, a Double-Tail latch allows lower power supply voltage range, with low-differential kickback. Otherwise, the Double-Tail exhibit a higher noise due to the integration in its first stage. Tested from -40$\degree$C to 200$\degree$C, the last stage of the proposed ADC topology does not need calibration up to 180$\degree$C. The encouraging results on this stage allows the re-use of the SAR to calibrate the previous stages. And considering the ADC, we estimate a possible resolution of 11.2-bits in 5 clock cycles per sample with an extension to 13.3-bits in 6 clock cycles with an estimated area of 0.12 $\rm mm^2$.

%-- paragraph sur les recommendations
The ADC test chip not being fabricated yet, a first step is the characterization of the ADC\@. From the results of the planned measurement session, the main goal is to push the architecture at higher sampling rates to then leverage the digital processing to enhance the sampling rate without changing the analog.
\end{multicols}
\end{small}
\end{mdframed}

\begin{textblock*}{161mm}(16mm,277mm)
\renewcommand{\familydefault}{\sfdefault}
\color{bordeau}
{\bf\noindent Université Paris-Saclay	         } \newline
\noindent Espace Technologique / Immeuble Discovery \newline
\noindent Route de l’Orme aux Merisiers RD 128 / 91190 Saint-Aubin, France 
\end{textblock*}

\begin{textblock*}{20mm}(182mm,268mm)
\includegraphics[width=20mm]{Figs/e.pdf}
\end{textblock*}